\chapter{Introduction}\label{ch:intro}

\section{Motivation}
A score specifies what is to be played, and a performance realises how it is
played. The difference between the two, in micro-timing, articulation and
dynamics, is where musical expression lives. Recovering that difference from
a recording is of scientific interest in its own right and of practical use
in analysis, in pedagogy, in expressive rendering, and as structured
supervision for audio models. Two commitments shape the whole thesis. First,
the score is treated as known, so the problem is the inverse of expressive
synthesis rather than full transcription. Second, uncertainty is a primary
output: a per-note estimate is only actionable together with a calibrated
error bar, and calibration is therefore evaluated throughout with the same
weight as accuracy.

The vocabulary of the thesis is a decomposition of the recording into four
layers. The score is what was written. The expression is how each written
note was played, a small vector of interpretable variables per note. The
timbre is what the instrument adds to a note once its expression is fixed,
and the noise is the rest. The score is given, the expression is the object
of inference, the timbre is a nuisance that the waveform setting integrates
out exactly rather than estimates, and the noise is modelled per note and
per channel. Every chapter can be placed in this picture. The piano and
strings-and-winds settings observe the expression layer through an
intermediate measurement, aligned MIDI in the one case and a pitch tracker
in the other, whereas the waveform setting observes the recording itself and
must separate all four layers at once.

Two statements of different kinds run through the thesis and are kept
distinct. The scientific objective is the representation itself: an
interpretable, score-aligned description of how a performance was played,
in per-note variables and later in within-note trajectories, that can be
inferred from a recording with explicit uncertainty. The methodological
contribution is the inference framework that recovers it, Bayesian
score-conditioned inverse synthesis, in which a score-structured prior is
combined with an observation model of the performance. The two are
complementary, and neither is asserted beyond what the experiments
validate. The representation is interpretable by construction but not
semantically validated, and the framework is validated stage by stage,
with its full integration at the waveform still open
(\S\ref{sec:future-phase3}).

The title names the object of study. With each note of the known score is
associated a small \emph{bundle} of expressive variables, timing,
articulation and dynamics at first and later intonation, vibrato and
loudness, so that a performance becomes a per-note field over the score
support and transcription becomes posterior inference over that field.
``Bundle'' is used in this everyday sense throughout. The geometric reading,
with per-note fibres and a connection over the score graph, is not developed
here.

\section{Problem statement}
Let a symbolic score fix, for each note, its pitch, its score onset and its
notated duration. A performance assigns to each note a small vector of
continuous expressive variables, here the onset residual $\tau$, the log
articulation ratio $\log r$ and the normalised velocity $v$. Given the score
and an aligned performance, the task is to infer the posterior over these
per-note variables, including for notes that are held out (masked), and to
do so with calibrated uncertainty. This is score-informed performance
transcription, or equivalently probabilistic inverse synthesis. In symbols,
with $S$ the score, $x$ the observation of the performance and $z$ the
latent expressive variables,
\begin{equation*}
  p(z \mid x, S) \;\propto\; p(x \mid z, S)\, p(z \mid S),
\end{equation*}
where $S$ is the score layer and $z$ the expression layer. The timbre and
noise layers live inside the observation model, in the waveform setting as
the harmonic amplitudes $a$, which are marginalised exactly, and the
residual $\varepsilon$. The score-structured prior $p(z \mid S)$ is the
subject of Chapter~\ref{ch:model}. The observation model $p(x \mid z, S)$ is
made progressively rawer over the thesis: per-note values read from aligned
MIDI (Chapter~\ref{ch:results}), tracker-derived per-note estimates with
their variances (Chapter~\ref{ch:phase2}), and the waveform itself through
a differentiable harmonic synthesizer (Chapter~\ref{app:phase3}). Inferring
the score itself, that is, full automatic music transcription, is out of
scope (\S\ref{sec:scope} and \S\ref{sec:future-amt}).

The thesis asks three questions, in the order in which the settings answer
them.
\begin{itemize}
  \item RQ1 (representation). What do interpretable per-note variables,
        and later within-note trajectories, capture of expressive
        performance, and what do they measurably discard?
  \item RQ2 (structured inference). Does a score-structured prior improve
        both the recovery of expressive variables and the calibration of
        their uncertainty, compared with independent or purely temporal
        baselines, when observations are missing (piano) or noisy (strings
        and winds)?
  \item RQ3 (direct waveform inference). Can a differentiable harmonic
        observation model recover expressive trajectories directly from
        the waveform, and can its uncertainty separate genuine performance
        ambiguity from the model's own mismatch?
\end{itemize}
RQ2 is the core of the thesis. Its hypotheses were fixed before any test
number was seen, and support was defined as improvement on both axes at
once, accuracy and calibration (\S\ref{sec:metrics}). RQ1 is answered by construction together
with a quantified account of what the representation discards, through the
within-note drift study of Appendix~\ref{app:drift} and the curve-level
extension of Chapter~\ref{app:phase3}. RQ3 is answered on the development
set only (Chapter~\ref{app:phase3}). Joining its observation model to the
RQ2 prior, which would close the inverse-synthesis loop, is the endpoint of
the programme rather than a result of this thesis. Its first span, the
score prior over per-note constants inside the waveform inference
(\S\ref{sec:phase3-priorvalue}), is studied on the development set over
three spans. There the value of the structure is the context's common
offset, with nothing measurable beyond it against that baseline, and a
significant advantage over a purely temporal prior at phrase scale. The
purely temporal baseline of RQ2 is also examined at whole-track scale on
both corpora, again on the development set. A naive note-order chain is
worse than the score graph, but the same graph with its pitch term removed
and its two-beat time window kept matches it
(Appendix~\ref{app:trackscale}), so the structure that the test-set results
reward is smoothness in score time rather than pitch proximity.

\section{The proposed model and the two-stage pipeline}\label{sec:twostage}
Two names recur throughout the thesis. \textbf{The proposed model} is the
thesis model: one multi-output Gaussian process over all notes and channels,
in which the score graph, the hand-built score features and the per-note
representation of the music model enter a single kernel, with every
hyperparameter learned from one per-piece marginal likelihood
(\S\ref{sec:gpfirst}). The music model is a small symbolic language model
trained from scratch before everything else. \textbf{The two-stage
pipeline} is the earlier development system that the proposed model
provably contains as a special case (\S\ref{sec:special-cases}). It
consists of a plug-in mean, a cross-piece point estimate of each note's
expression, followed by a separate per-channel graph Gaussian process on
the residual.

\section{Contributions}
\begin{enumerate}
  \item A multi-output graph Gaussian process for expressive performance
        (\S\ref{sec:gpfirst}). Channels are coupled by coregionalization,
        the graph kernel is a spectral function of the score-graph
        Laplacian, and all side information, hand-built score features and
        learned music-model embeddings alike, enters as marginalized
        Bayesian linear means. Every hyperparameter is learned by the exact
        per-piece evidence, and the posterior and its calibration follow in
        closed form. Its hypotheses were fixed in advance and evaluated
        once on test pieces untouched by any modelling decision.
  \item An attribution of the model's value to its parts
        (\S\ref{sec:attribution}). Per-piece Bayesian weighting of
        informative features does the recovering, the graph makes the
        confidence calibrated, with its calibration contribution supported
        independently on the test set, and the music model trained from
        scratch earns its place as a feature provider inside the kernel.
        Each of these statements is isolated by nested ablations, including
        the entire two-stage pipeline as a provable special case. The graph
        term is attributed one step further on the development set of both
        corpora: the same graph without its pitch term reproduces its value
        (Appendix~\ref{app:trackscale}), so what the graph contributes is
        smoothness in score time over about two beats.
  \item A map of where the structure helps, obtained from six downstream
        tasks with objective metrics and re-validated under the final
        model. It yields two boundaries (\S\ref{sec:boundary}): structure
        helps note-level judgements but not whole-performance summaries,
        and per-piece adaptation helps interpolation but not extrapolation
        from an excerpt.
  \item An evaluation protocol in which leakage control is itself a result
        (\S\ref{sec:hygiene}). An evaluation leak in the music-model
        read-out was found and corrected, and structural tests together
        with bitwise end-to-end audits now exclude it. The protocol further
        comprises contamination filtering between corpora, robustness
        safeguards that validate themselves, and a development/test split
        that exposed the inflation of numbers reused for selection.
\end{enumerate}
The waveform endpoint, the full inverse-synthesis loop with the score-graph
prior inside the waveform inference, is the intended destination of the
programme and is not presented as a contribution. Its components are
studied on the development set (Chapter~\ref{app:phase3}), and the
integration is built only in its first span, where the score-graph prior
couples the per-note intonation constants inside the waveform inference on
short chains. In that span, with hypotheses stated in advance and after a
flaw in the first design had been found and corrected, the value of the
prior is the transfer of the track's shared tuning offset to a note whose
audio is absent. Beyond that common mode, score structure adds nothing
measurable for intonation at chain or phrase scale or for timing at chain
scale, although at phrase scale it is better than a purely temporal prior.
These are three development studies
(\S\ref{sec:phase3-priorvalue}). The curve-level coupling, the remaining
channels and whole-performance scale remain open: the within-note curves
still carry independent priors, and the graph meets the full waveform
posterior only post hoc, as an observation channel of the
Chapter~\ref{ch:phase2} bundle. A geometric ``connection-Laplacian''
generalisation of the graph prior likewise remains future work
(Chapter~\ref{ch:future}).

\section{Scope and non-goals}\label{sec:scope}
The core result is scoped to piano and to symbolic performance variables
read from MIDI, namely onset timing $\tau$, articulation $\log r$ and
dynamics $v$. The score support is known, that is, the note identities,
the score onsets and the score durations are given. Evaluation is on
held-out notes and pieces and always includes calibration. Four things are
explicitly outside the scope of this first version.
\begin{itemize}[noitemsep]
  \item discovering unknown score support from audio, which is full
        automatic music transcription and a future extension
        (\S\ref{sec:future-amt}).
  \item performer-controlled intonation or timbre \emph{from piano}, which
        do not exist on a fixed-pitch keyboard.
  \item arbitrary instruments and room acoustics.
  \item any suggestion that the model performs general automatic music
        transcription.
\end{itemize}
Audio input and output and continuously pitched instruments enter only
with the strings-and-winds and waveform settings. The work is staged so
that each stage produces a defensible result before the next difficulty is
taken on, and the stages are ordered toward the waveform endpoint: each
isolates one new difficulty on the way to inference from sound alone.
Table~\ref{tab:phases} records the stages and their status. The thesis
names the stages by their content.\footnote{In the code, the results files
and the names of archived artifacts the four stages are numbered 0--3 in
this order: the music model, piano, strings and winds, the waveform.}

\begin{table}[h]\centering
\caption{The four stages and their status.}
\label{tab:phases}
\footnotesize
\begin{tabular}{@{}p{2.6cm}p{3.8cm}p{3.7cm}p{3.4cm}@{}}
\toprule
Stage & Scope & Adds & Status \\
\midrule
music model (foundation) & From-scratch symbolic music model (backbone + representations) &
Per-note embeddings $h_i$ (a feature kernel of the GP; the plug-in mean $\muLM$ of the two-stage baseline) & Implemented \\
piano (stage 1): isolating the prior & Timing, articulation, dynamics from aligned score+performance &
Score-graph GMRF prior, closed-form posterior, calibrated uncertainty &
Evaluated on a held-out test set (Chapter~\ref{ch:results}) \\
strings and winds (stage 2): measurement uncertainty & Intonation, vibrato, loudness on continuously pitched instruments &
$f_0$-derived channels ($c_i$ cents, vibrato, $\ell_i$) $+$ $\tau$, $d^{\mathrm{vib}}$ & Evaluated on a held-out test set (URMP): intonation recovery, vibrato calibration and coverage supported, the secondary timing hypothesis not supported (Chapter~\ref{ch:phase2}) \\
waveform (stage 3): the endpoint & Differentiable synthesizer as audio likelihood &
Harmonic amplitudes $a$, audio $x$ & Development studies: scalar intonation, whole pitch curves, joint note boundaries (Chapter~\ref{app:phase3}); no held-out evaluation yet \\
\bottomrule
\end{tabular}
\end{table}

\section{Thesis outline}
Chapter~\ref{ch:related} positions the work against expressive-performance
modelling, alignment, structured priors, symbolic music models and
differentiable synthesis. Chapter~\ref{ch:model} gives the model and its
derivation in the piano setting. Chapter~\ref{ch:data} describes the data
and the evaluation protocol, including contamination control.
Chapter~\ref{ch:results} reports the piano result and the methodological
findings, and Chapter~\ref{ch:downstream} maps the downstream reach.
Chapter~\ref{ch:phase2} carries the same prior, unchanged, to real
recordings and reports the strings-and-winds results on the held-out test
set. Chapters~\ref{ch:discussion}--\ref{ch:conclusion} discuss, situate and
conclude. The appendices hold the notation summary (\ref{app:notation}),
the per-channel form of the two-stage lineage (\ref{app:lineage}), the
relocated piano and strings-and-winds development studies
(\ref{app:phase1dev}, \ref{app:phase2dev}) and the calibration profile
(\ref{app:extra}).

% =====================================================================
